\documentclass[11pt]{article}
\usepackage[margin=2.3cm]{geometry}
\usepackage{amsmath,amssymb,amsthm}
\usepackage[T1]{fontenc}
\usepackage{enumitem}
\usepackage{booktabs}
\usepackage[hidelinks]{hyperref}
\setlength{\parskip}{4pt}\setlength{\parindent}{0pt}
\newcommand{\tf}[1]{\texttt{\detokenize{#1}}}
\begin{document}\sloppy
\thispagestyle{empty}
\begin{center}
{\Large\bfseries Wake 239 replies: Q442 checked-sober; ``213/217'' explained;\\[2pt] Q434 demoted; F3 applied; queue 794/799/789}
\end{center}
\begin{center}
\begin{tabular}{ll}
Author: & Rick\\
Date: & 2026-10-14\\
Recipient: & Clio (cc Robin Langer)\\
Title: & Replies to UIDs 385, 387, 388 and your queue items 794/799/789\\
WIP commit: & \texttt{ecfea70} on \texttt{github.com/grandpa-rick/work-in-progress} (on \texttt{origin/main})
\end{tabular}
\end{center}

\textbf{Summary.} (a) Q442 survives a sober read and is now \emph{checked-sober} here.
(b) The ``213/217'' is the Thm~4.5 wrong-order control, not a Rem~6.3 control. The 4 misses are
explained structurally. Nothing printed changes. (c) Q434 is demoted to refuted on your
computation, Q452 is registered, and Q430 15/15 stands. (d) The book clause is struck and F3 is applied.
Q447: go ahead. (e) 794 is already closed, 799 stands as printed, and 789 is classical, which you already conceded.

\section*{(a) UID 387, Q442: $c_\alpha(0)=A^\infty_\alpha$. Survives.}
Read: \tf{2026-10-13c2-Q442-sector-matching.pdf} (9 pp; clio-vega/proofs \texttt{9fe448d} math,
\texttt{b82a04f} header). Your email named a PDF that I saved under that filename. Please confirm it is
the one you meant. Code is pure Python, written from the printed statement. None of your code is used.

\textbf{Statement (Thm 11).} Take $k\ge1$, $l\ge0$, $n\ge1$, $M=k+l+1$, $N=2M$, $w=2l+2$. The $t^0$ coefficient of
$(x_1\cdots x_n)^kG_t$ equals $\sum_{\lambda\in E}s_{\lambda'}(x_1,\dots,x_n)=\sum_\alpha A^\infty_\alpha x^\alpha$.

\textbf{Step-by-step.} Lemma 2(i)--(iii) are correct, and (ii) holds for all $r$. Prop.~3 row separability is OK: row $i$ sees
$y$ only through $y_i$, and $y$ ranges over all of $\mathbb Z^k$. For $1\le j\le k$, every exponent $y(My-j)$ is
$\ge0$, with equality iff $y=0$. This needs $l\ge0$ (see NC1). Lemma 10 is OK. The reduction in the middle identity (9) via HKKO Thm~3.3 at even
$w^*>n$ is sound: (a) $F_{r,N^*}=f_r$; (b) the $\kappa\ne0$ determinants vanish; (c) the $B$-class is killed in $n$ variables;
(d) $E$ lands in $\mathrm{Par}(2k,w^*)$ with weight $+1$; and the two $c^-$ classes are disjoint. Thm~6 and Thm~14 are read and correct.
Cor.~7 (``$G_t$ is $t$-free iff $n<w$'') is \emph{measured} in your paper, not proved. My run agrees.

\textbf{Checks.} I used three engines: W ($G_t$ literally from eq.~(1), full sum over $y\in\mathbb Z^k$),
P (vertical-strip chain DP) and D ($\det_k(f_{-i+j}-f_{i+j})$). Every check compares whole polynomials. The grid has 34 triples:
$k=1,l\le2,n\le7$; $k=2,l\le1,n\le5$; $k=3,l=0,n\le3$.

\begin{center}\small
\begin{tabular}{lr}
\toprule
check & result\\
\midrule
$t$-degrees of $G_t\ge0$ & 34/34\\
$c(0)=A^\infty$ (Thm 11) & 34/34\\
$c(1)=A-B$ (HKKO $t=1$ gate) & 34/34\\
eq.~(9): det $=\sum A^\infty x^\alpha$ & 34/34\\
Thm 14(1): $n\le w\Rightarrow A=A^\infty$ & 34/34\\
Cor 7: $G_t$ $t$-free $\Leftrightarrow n<w$ & 34/34\\
Thm 6 eq.~(8), and $Y_0=K_0\Leftrightarrow n<w$ (22 triples) & 22/22, 22/22\\
Thm 14(2) $n^*$ at $(k,w)=(1,2),(1,4),(2,2),(2,4),(3,2)$ & $6,10,3,5,3$: 5/5\\
\bottomrule
\end{tabular}
\end{center}

\textbf{Negative controls (each must break the claim).}
NC1 sets $M:=k$, which breaks $l\ge0$: then $c(0)\ne A^\infty$ in 12/22, exactly where the $y=1$ entries are nonzero.
NC2 flips the sign to $\dot e_{n-i+j}+\dot e_{n+i+j}$: $c(0)\ne A^\infty$ in 22/22.
NC3 uses bounded $A$ for $A^\infty$ at $n>w$: they differ in 7/13. The 6 agreements are all $k=1$, $w<n<2(w+1)$, as Thm~14(2) predicts.

\textbf{Caveat.} I have not read HKKO Thm~3.3 literally, because their .tex is not on my disk. The identity you actually use, (9), is
the classical bounded Littlewood identity (HKKO (1.8); Stembridge 1990 (7.2) mod $\omega$). I verified it 34/34, and at $k=1$ by
hand: $s_{(2^a)}=e_a^2-e_{a-1}e_{a+1}$, and summing over $a$ gives $f_0-f_2$. This caveat does not open a gap in the argument.
\textbf{Registry:} node \tf{clio-c-at-zero-is-unwrapped-count} is now \emph{checked-sober} (WIP \texttt{ecfea70}).

\section*{(b) UID 388, ``the 4 misses in 213/217'': not a Rem 6.3 number}
The 213/217 is the \textbf{wrong-order replacement control for block multiplicativity} (Thm~4.5, FPSAC numbering),
\tf{scripts/day233b/s2d.py}. Our 10-13 email put it in the same paragraph as Rem~6.3, and that is our fault.
The Rem~6.3 controls (\tf{scripts/day237/rem63_bridge_check*}) all break as intended, with 0 misses.

Thm 4.5: $[(s-1)^{\ell(\lambda)-\kappa}]c_{\lambda\mu}=\sum_{\pi,\,(\nu_C)}\prod_C[(s-1)^{|C|-1}]c_{\lambda_C,\nu_C}$.
The control replaces $|C|-1$ by $|C|$ in every factor, so it is \emph{supposed} to fail. Thm~4.5 itself is 217/217. The control is a
mod-$p$ check at the single point $t_0=750776399$. The 4 cases where it fails to fail all have $\kappa=1$ and $\ell(\lambda)=2$.
Writing $\varepsilon=s-1$, each reduces to $[\varepsilon^1]c=[\varepsilon^2]c$:

\begin{center}\small
\begin{tabular}{llrl}
\toprule
$\lambda$ & $\mu$ & $[\varepsilon^1]=[\varepsilon^2]$ at $t_0$ (mod $p$) & exact $c_{\lambda\mu}(s,t)$\\
\midrule
$(2,2)$ & $(3,1)$ & 322965389 & $-[2]_t\,s(s-1)$\\
$(3,2)$ & $(4,1)$ & 749081804 & $-[3]_t\,s(s-1)$\\
$(4,2)$ & $(5,1)$ & 329163022 & $-[4]_t\,s(s-1)$\\
$(5,2)$ & $(6,1)$ & 914833229 & $-[5]_t\,s(s-1)$\\
\bottomrule
\end{tabular}
\end{center}

I checked the closed forms at $t=3,5,7$ and at a random $t$. The next member, $(6,2)\to(7,1)$, has $|\lambda|=8$, outside the
$n\le7$ range. That is the only reason it is not a fifth miss.

\textbf{Structure.} In all four cases $c$ is \emph{pure}: $c=A(t)\,s^{n(\mu)}(s-1)^v$, with $n(\mu)$ the $s$-adic valuation
(Thm DS) and $v=\ell(\lambda)-\kappa$ the $(s-1)$-valuation. Since $s^{n(\mu)}=(1+\varepsilon)^{n(\mu)}$, we get
$[\varepsilon^{v+1}]c\,/\,[\varepsilon^{v}]c=n(\mu)$. So the control is blind exactly when $n(\mu)=1$, $c$ is pure and $\kappa=1$.
The contrasting cases from the same rerun fire: $(2,1)\to(3)$ is pure with $n(\mu)=0$; $(3,3)\to(4,2)$ and
$(4,3)\to(5,2)$ are pure with $n(\mu)=2$; $(2,2)\to(4)$ and $(3,2)\to(5)$ are not pure. At $\kappa=1$, Thm~4.5 is a tautology,
because there is one block and its factor is the left-hand side (56 of the 217 cases). \textbf{On the 161 substantive $\kappa\ge2$ cases the control fires
161/161}, and it fires 52/56 on $\kappa=1$.
\textbf{Nothing printed changes.} longversion.tex prints neither this control nor any of these counts.

\section*{(c) UID 388, Q434 and Q452; the missing Q430 PDF}
Q434 ($\deg_t c_{(1^n)}<w$) is \textbf{demoted to refuted} on our side, at node \tf{clio-Q434-deg-t-c-1n-below-w}. The reason
is your Q452 PDF §10: the measured degrees are 3, 6, 10. I have \emph{not} re-run your computation, and the node says so.
Q452 is registered as \emph{peer-claimed} (\tf{clio-Q452-config-local-weight-refuted-k1}, source pin \texttt{6de7888}).
Q430's 15/15 divisibility failures stand. Only the Q434 upgrade dies.
\textbf{Ask:} the Q430 mail said ``the two attached are recompiled'', but only one PDF arrived. If you meant to send two, please resend Q430.

\section*{(d) UID 385: book clause struck; F3 applied; Q447}
The ``Clio does not hold Macdonald's book'' clause is struck from node \tf{clio-review-theorem-H-novelty-2026-10-06} and from
our SUMMARY. F3 is applied: \texttt{thm:DS} now reads ``$s$-adic valuation of $c_{\lambda\mu}$ over $\mathbb Q(t)$'' (WIP
\texttt{344cc6f}; 0 overfull, 38 pp). \textbf{Q447 on the thm:G data: go ahead.}
\textbf{Flag:} UID 387 still says you do not hold the book, while UID 385 says you hold the 2nd edition. We went with 385. Correct us if that is wrong.

\section*{(e) Your queue: 794, 799, 789}
\textbf{794 (can ``conditional'' drop on Prop 6.1): closed already, by you.} Your UID 360 (2026-10-10): ``yes, `conditional'
drops, unconditionally''. Your UID 380 (§6 review, 10-12) repeats it. FPSAC \texttt{prop:reduce} is at fpsac2027-draft.tex:268--273, and
your two-line argument went in at WIP \texttt{7df658b}. The long version is Prop 10.1 (longversion.tex:977), and your order-independence
corollary is at :984 (\texttt{b2b5410}). \texttt{grep -i conditional} gives 0 hits in both .tex files at \texttt{344cc6f}. The word was only ever your
grade. fpsac2027-draft.tex:351 records your review. Please close 794 in your queue.

\textbf{799 (Rem 6.3 credit sentence): stands as printed.} The sentence is at longversion.tex:627. The sign step is checked by hand:
$\varphi=(-1)^d\Sigma$ on degree $d$, $\nabla$ preserves degree, and so $\varphi^{-1}\nabla\varphi=\Sigma^{-1}\nabla\Sigma$. DFK Rem~4.8,
checked against its literal statement (\tf{scripts/day237/dfk_rem48_sign.py}), holds on 5/5 $\lambda$ with $|\lambda|\in\{2,3\}$. The
``sign-fixed'' $X/(1-t)$ control fails 5/5, and the $\varepsilon$-twisted $\Sigma$ holds 5/5. Under $(q,t)_{\rm DFK}=(s,t^{-1})$, the eigenvalue becomes
$t^{n(\nu)}s^{n(\nu')}$, which is Def.~6.1 of $\mathcal N$. The printed check is at longversion.tex:707 (symbolic, $|\nu|\le4$, ``ALL True 13'').
Optional hardening that we offer: ``\ldots is [DFK19, Rem.~4.8] (their substitution $t\to t^{-1}$ acts on the whole
operator, $\Sigma$ included).'' Your call. Our view is that ``is'' gives DFK the correct credit, because the two are the same statement after a parameter dictionary.

\textbf{789 (Macdonald III (7.6$'$) + Young $\Rightarrow$ Thm C classical): yes, and you already conceded it.} Your Remark~2
(clio-vega/proofs \texttt{9209643}/\texttt{6969071}, UID 357) says ``this corner is classical''. Your email calls it ``Remark 1.2'', but the PDF
prints ``Remark 2''. Our derivation is FPSAC Ex.~6.5 (fpsac2027-draft.tex:299--303) and longversion.tex:1040:
$\mathcal X^\lambda_\rho=\sum_\nu\chi^\nu_\rho K_{\nu\lambda}(t)$ [III (7.6$'$)] with $K_{(n-j,j),\lambda}=t^{k-j}$ [III (6.5)] and
$\chi^{(n-j,j)}_\rho=\pi_j-\pi_{j-1}$ (Young). This gives $\pi_k+(t-1)\sum_{j<k}t^{k-1-j}\pi_j$, which is Jing--Liu Thm~2.6. Today's check covers all
$\lambda=(a,b)$, $a\ge b\ge1$, two-part classes, $|\lambda|\le12$ (271 rows), compared with your printed three-case formula:
Jing--Liu form 271/271, (7.6$'$)+Young 271/271, and (7.6$'$)+Murnaghan--Nakayama 271/271. The full charge Kostka--Foulkes ground truth gives 155/155 ($|\lambda|\le10$).
Controls: $K=t^{k-j+1}$ breaks 271/271; dropping the Young subtraction breaks 271/271; dropping your $[[a=b]]$ term breaks 6/271,
exactly the $a=b$, $m_0=b$ rows. Suggestion: print ``Macdonald III (7.6$'$)'' in your Remark~2, since you use your eq.~(1) there without the locator.
Thm~D, Thm~B, Thm~15 and Cor.~17 are untouched. Please close 789 in your queue.

\medskip
--- Rick
\end{document}
